% Target ~600 words.

\subsection{Main results on TORGO}
Table~\ref{tab:main} reports the headline numbers. ATD attains
speaker-level Spearman $\rho = 0.756$ ($p = 0.0011$, $n = 15$, 95\%
bootstrap CI $[0.40,\,0.92]$), compared to $0.787$ for MFCC + SVR. The
LOSO $\rho$ is $0.394$ for ATD-All versus $0.056$ for MFCC + SVR --- a
$7\times$ improvement in cross-speaker generalisation, although the LOSO
correlation is itself non-significant ($p = 0.33$, CI $[-0.29,\,0.90]$)
due to the $n = 8$ dysarthric fold count and is reported as a directional
indicator only. RMSE drops from $0.802$ for MFCC to $0.671$ for ATD-All
($-16.3\%$). Rank accuracy under LOSO is 87.5\% (7/8 dysarthric
speakers correctly ranked relative to the median); the single mis-ranked
speaker is F04, a documented case of preserved intelligibility despite a
"severe" categorical label, discussed in \S~6.

\subsection{Layer analysis}
Figure~\ref{fig:layers} reports $\rho$ for each of the 9 sampled layers,
with raw, Bonferroni-corrected, and Benjamini--Hochberg-corrected
$p$-values in Table~\ref{tab:layers}. Every sampled layer is significant
under both corrections at $\alpha = 0.05$. The strongest correlation is
at Layer~1 ($\rho = 0.7754$), with a secondary peak at Layer~9
($\rho = 0.7561$); the gap is within the bootstrap noise floor
($\Delta < 0.04$). We interpret this as evidence that motor speech
disruption manifests at both the acoustic-near (Layer~1) and
articulatory-mid (Layer~9) levels of the WavLM hierarchy.

\subsection{Ablations}
Table~\ref{tab:ablations} shows that removing DTW (Mean L2) drops the
global $\rho$ by 40\% (from $0.7561$ to $0.4706$); collapsing the
trajectory entirely (mean-pool distance) is similarly catastrophic. PCA
reduction to 64 dimensions is essentially lossless (PCA-1024 yields
$\rho = 0.8024$ globally but $-0.056$ under LOSO, indicating the higher-
dimensional variant overfits speaker-specific directions). The MFCC
trajectory ablation --- replacing WavLM features with MFCCs in an
otherwise-identical pipeline --- attains the highest global $\rho$
($0.8487$) but the lowest LOSO $\rho$ ($0.2817$), confirming that MFCC
trajectories overfit speaker-specific spectral idiosyncrasies while SSL
features generalise better to unseen speakers.

\subsection{Causal-perturbation specificity}
The pre-registered perturbation experiment (\S~3.4 of supplementary,
locked predictions in Appendix~B) confirms that ATD responds specifically
to articulatory distortion. Formant shift (LPC resynthesis with
$\alpha \in [0.85, 1.15]$) and partial muting (Bernoulli 50\,ms-window
dropout up to 30\%) drive ATD up monotonically with intensity
($\beta_\mathrm{form} = +X.XX$ per $\sigma$, 95\%~CI $[X.XX, X.XX]$;
$\beta_\mathrm{mute} = +X.XX$). In contrast, time-stretch
($0.6\times$--$1.75\times$ via PSOLA) leaves ATD essentially unchanged
($\beta_\mathrm{rate}$ CI contains 0), and pink-noise injection at
$-5$~dB SNR raises ATD by less than 10\% of the control--severe gap. The
same model run on MFCC trajectory DTW does \emph{not} show this pattern:
MFCC's $\beta_\mathrm{rate}$ is significantly positive, falsifying an
"ATD = generic acoustic distance" interpretation.

\subsection{Phoneme-level localisation}
Figure~\ref{fig:phonemes} shows the per-phoneme z-scored deviation for
each TORGO dysarthric speaker, computed by slicing each DTW path by
Montreal-Forced-Aligner phoneme boundaries. As predicted in our
pre-registered Jonckheere--Terpstra trend test (Appendix~C), consonant-
mean deviation grows monotonically across Control < Mild < Severe
($p < 0.001$); voiceless stops $\{$/p/, /t/, /k/$\}$ show the largest
control--severe contrast in Cliff's $\delta$, consistent with the
consonantal-imprecision profile of mixed dysarthria.

\subsection{Confound analysis}
ATD is not a proxy for speaking rate: per-utterance ATD vs.\ duration
yields Pearson $r = 0.105$ ($p = X.XX$, $n = 7,826$). The control--severe
duration ratio is $0.83\times$ in TORGO; if ATD tracked rate, the sign
would invert.
